%% kapa: exam=1 round=24 year=2013 period=1 form=A subject=civil_law subject_ko=민법 seq=9 number=9
%% source: Q-Net 감정평가사 기출문제 (cst00309) artlSeq=5073802 file=extracted/1008/5073802/2013년 제24회 감정평가사 1차시험 기출문제/1교시(A형).hwp
%% conversion: kordoc 4.24.51 -> tools/gen1.py (auto); review=proofread-sonnet; flags=none
\begin{question}{9}
법률행위의 무효와 취소에 관한 설명으로 옳지 않은 것은? (다툼이 있으면 판례에 의함)
\begin{choices}
  \item 토지거래허가를 배제하거나 잠탈하는 내용의 토지매매계약은 확정적으로 무효이다.
  \item 무효인 법률행위에 의한 법률효과를 침해하는 것처럼 보이는 위법행위가 있더라도 그 법률효과의 침해에 따른 손해가 인정되지 않는다.
  \item 무효인 재산상 법률행위를 당사자가 무효임을 알고 추인한 경우 제3자에 대한 관계에서도 처음부터 유효한 법률행위가 된다.
  \item 무능력을 이유로 법률행위가 취소된 경우 악의의 무능력자도 그 행위로 인하여 받은 이익이 현존하는 한도에서 반환할 책임이 있다.
  \item 강박에 의한 의사표시가 적법하게 취소된 후에 그 의사표시를 추인하는 경우 그 취소원인이 소멸한 후에 하여야만 추인의 효력이 있다.
\end{choices}
\end{question}
\answer{3}
